\chapter{Calcul intégral}

\noindent\textit{Intégrer, c'est « sommer » une infinité de contributions infinitésimales :
l'intégrale d'une fonction positive mesure une \textbf{aire}. Le lien fondamental avec les
primitives — l'intégrale de $a$ à $b$ vaut $F(b)-F(a)$ — fait du calcul intégral un outil
de mesure (aires, volumes, valeurs moyennes) et un terrain privilégié des problèmes de Baccalauréat
de la série C. Ce chapitre rassemble tout l'outillage : définition de l'intégrale, théorème
fondamental, propriétés (linéarité, Chasles, positivité, inégalité de la moyenne), valeur moyenne,
intégration par parties, suites d'intégrales, aires et volumes de révolution.}
\vspace{8pt}

% ═══════════════════════════════════════════════════════════════
\section{Primitives : rappels et compléments}

Le calcul intégral repose entièrement sur la notion de \textbf{primitive}. On la rappelle ici
avant d'en faire l'usage central.

\begin{definition}[Primitive d'une fonction]
Soit $f$ une fonction définie sur un intervalle $I$. Une fonction $F$ dérivable sur $I$ est une
\textbf{primitive} de $f$ sur $I$ si $F'=f$ sur $I$.
\end{definition}

\begin{theoreme}[Ensemble des primitives]
Si $f$ admet une primitive $F$ sur un intervalle $I$, alors $f$ admet une \textbf{infinité} de
primitives sur $I$ : ce sont exactement les fonctions $x\mapsto F(x)+k$, $k\in\R$. De plus,
pour $x_0\in I$ et $y_0\in\R$ donnés, il existe une \emph{unique} primitive $G$ telle que $G(x_0)=y_0$.
\end{theoreme}

\begin{remarque}
Toute fonction \textbf{continue} sur un intervalle $I$ y admet des primitives (résultat admis,
conséquence du théorème fondamental ci-dessous). C'est pourquoi, dans tout ce chapitre, on
intègre des fonctions continues.
\end{remarque}

\begin{propriete}[Primitives usuelles]
On note $F$ une primitive de $f$ sur un intervalle où $f$ est définie ($k$ désigne une constante).
\begin{center}\small
\begin{tabular}{@{}llll@{}}
\toprule
$f(x)$ & $F(x)$ & $f(x)$ & $F(x)$\\
\midrule
$x^{n}\ (n\ne-1)$ & $\dfrac{x^{n+1}}{n+1}$ & $\dfrac1x\ (x>0)$ & $\ln x$\\[5pt]
$e^{x}$ & $e^{x}$ & $\dfrac1{x^{2}}$ & $-\dfrac1x$\\[5pt]
$\sin x$ & $-\cos x$ & $\cos x$ & $\sin x$\\[5pt]
$\dfrac1{\sqrt x}$ & $2\sqrt x$ & $\dfrac1{\cos^{2}x}=1+\tan^2 x$ & $\tan x$\\
\bottomrule
\end{tabular}
\end{center}
\end{propriete}

\begin{propriete}[Primitives « avec $u$ » (formes composées)]
Soit $u$ une fonction dérivable. Alors une primitive de :
\[
u'u^{n}\ (n\ne-1)\ \text{est}\ \frac{u^{n+1}}{n+1};\quad
\frac{u'}{u}\ (u>0)\ \text{est}\ \ln u;\quad
u'e^{u}\ \text{est}\ e^{u};
\]
\[
\frac{u'}{\sqrt u}\ \text{est}\ 2\sqrt u;\quad
u'\cos u\ \text{est}\ \sin u;\quad
u'\sin u\ \text{est}\ -\cos u.
\]
\end{propriete}

\begin{exemple}
Une primitive de $f(x)=\dfrac{2x}{x^{2}+1}$ est $\ln(x^{2}+1)$ (forme $u'/u$).
Une primitive de $g(x)=x\,e^{x^{2}}$ est $\tfrac12 e^{x^{2}}$ (forme $u'e^{u}$ à un facteur $\tfrac12$ près,
car $(x^2)'=2x$).
\end{exemple}

\begin{aretenir}
\textbf{Réflexe « avec $u$ »} : devant une intégrale, on cherche d'abord à reconnaître l'une des
formes $u'u^n$, $\frac{u'}{u}$, $u'e^u$, $u'\cos u$\dots\ Cela évite des intégrations par parties
inutiles. On ajuste seulement un \emph{facteur constant} si nécessaire.
\end{aretenir}

% ═══════════════════════════════════════════════════════════════
\section{Intégrale d'une fonction continue}

\begin{definition}[Intégrale d'une fonction continue]
Soit $f$ une fonction \textbf{continue} sur un intervalle $I$ et $a,b\in I$. Si $F$ est
\emph{une} primitive de $f$ sur $I$, le nombre $F(b)-F(a)$ ne dépend pas de la primitive choisie.
On l'appelle \textbf{intégrale de $f$ de $a$ à $b$} :
\[
\int_a^b f(x)\dd x=\big[F(x)\big]_a^b=F(b)-F(a).
\]
$x$ est une variable \emph{muette} : $\displaystyle\int_a^b f(x)\dd x=\int_a^b f(t)\dd t$.
Les nombres $a$ et $b$ sont les \textbf{bornes} de l'intégrale.
\end{definition}

\begin{remarque}
Le résultat est indépendant de la primitive : si $G=F+k$, alors $G(b)-G(a)=(F(b)+k)-(F(a)+k)=F(b)-F(a)$.
La constante $k$ « disparaît » dans le calcul d'une intégrale.
\end{remarque}

\begin{exemple}
$\displaystyle\int_1^2 x^2\dd x=\left[\frac{x^3}{3}\right]_1^2=\frac{8}{3}-\frac13=\frac{7}{3}$.\;
$\displaystyle\int_0^{\pi}\sin x\dd x=\big[-\cos x\big]_0^{\pi}=-(-1)-(-1)=2$.
\end{exemple}

\begin{theoreme}[Théorème fondamental du calcul intégral]
Soit $f$ continue sur $I$ et $a\in I$. La fonction
\[
\Phi:x\longmapsto\int_a^x f(t)\dd t
\]
est l'\textbf{unique} primitive de $f$ sur $I$ qui s'annule en $a$ : elle est dérivable et
$\Phi'(x)=f(x)$ pour tout $x\in I$, avec $\Phi(a)=0$.
\end{theoreme}

\begin{remarque}
Ce théorème est le cœur du chapitre : il relie \emph{dérivation} et \emph{intégration}, deux
opérations réciproques. Il justifie aussi la formule $\int_a^b f=F(b)-F(a)$ : si $\Phi(x)=\int_a^x f$,
alors $\Phi$ est une primitive de $f$, donc $\int_a^b f=\Phi(b)=\Phi(b)-\Phi(a)=F(b)-F(a)$ pour
toute primitive $F$.
\end{remarque}

\begin{exemple}
Soit $\Phi(x)=\displaystyle\int_1^x \dfrac{\dd t}{t}$ pour $x>0$. D'après le théorème fondamental,
$\Phi'(x)=\dfrac1x$ et $\Phi(1)=0$ : on reconnaît $\Phi(x)=\ln x$. De même
$\dfrac{\dd}{\dd x}\displaystyle\int_0^x e^{-t^2}\dd t=e^{-x^2}$, bien qu'on ne sache pas exprimer
cette intégrale avec les fonctions usuelles.
\end{exemple}

\begin{theoreme}[Interprétation comme aire]
Si $f$ est continue et \textbf{positive} sur $[a,b]$, alors $\displaystyle\int_a^b f(x)\dd x$ est l'\textbf{aire}
(en unités d'aire) du domaine compris entre la courbe $\mathcal C_f$, l'axe des abscisses et les
droites $x=a$ et $x=b$.
\end{theoreme}

\begin{illus}
\begin{tikzpicture}[>=Stealth]
\begin{axis}[width=8.4cm,height=5cm,axis lines=middle,xlabel={\small $x$},ylabel={\small $y$},
  xmin=-0.4,xmax=4,ymin=-0.4,ymax=3.4,xtick={1,3},xticklabels={$a$,$b$},ytick=\empty,
  clip=false]
\addplot[onbo,line width=1pt,domain=0:3.7,samples=80]{0.3*x^2-0.1*x+0.8};
\addplot[onbo!35,domain=1:3,samples=60,fill=onbsoft]{0.3*x^2-0.1*x+0.8}\closedcycle;
\addplot[onbo,line width=1pt,domain=1:3,samples=60]{0.3*x^2-0.1*x+0.8};
\node[onbo,anchor=west] at (axis cs:3.1,2.6){\small $\mathcal C_f$};
\node[onbink] at (axis cs:2,1.0){\small $\displaystyle\int_a^b\! f$};
\end{axis}
\node[align=left,text width=3.3cm,anchor=west] at (8.6,1.6)
{\small L'aire coloriée sous la courbe, entre $x=a$ et $x=b$, vaut $\int_a^b f(x)\dd x$.};
\end{tikzpicture}
\fig{Figure 1 — Aire sous une courbe positive : interprétation de l'intégrale.}
\end{illus}

\begin{remarque}
Sur un repère orthogonal d'unités $i$ (sur $Ox$) et $j$ (sur $Oy$), l'\textbf{unité d'aire} (u.a.)
vaut $i\times j$. Une intégrale donne un nombre de u.a. ; on multiplie par la valeur d'une u.a.
(en cm$^2$) pour une aire « réelle ».
\end{remarque}

\begin{remarque}[Approche par les sommes de Riemann]
L'aire sous $\mathcal C_f$ s'obtient comme \emph{limite} de la somme des aires de fins rectangles :
en découpant $[a,b]$ en $n$ sous-intervalles de largeur $\frac{b-a}{n}$,
\[
\int_a^b f(x)\dd x=\lim_{n\to+\infty}\frac{b-a}{n}\sum_{k=0}^{n-1} f\!\left(a+k\,\frac{b-a}{n}\right).
\]
Cette idée — « sommer une infinité de tranches infiniment fines » — est l'origine du symbole
$\int$ (un « S » allongé, pour \emph{somme}).
\end{remarque}

\begin{illus}
\begin{tikzpicture}[>=Stealth]
\begin{axis}[width=8.4cm,height=5cm,axis lines=middle,xlabel={\small $x$},ylabel={\small $y$},
  xmin=-0.3,xmax=4.2,ymin=-0.3,ymax=3.4,xtick=\empty,ytick=\empty,clip=false]
% rectangles de Riemann (méthode des rectangles à gauche)
\addplot[ybar interval,fill=onbsoft,draw=onbo!55,line width=0.4pt,domain=0.3:3.5,samples=9]{0.3*x^2-0.1*x+0.7};
\addplot[onbo,line width=1.1pt,domain=0:4,samples=90]{0.3*x^2-0.1*x+0.7};
\node[onbo,anchor=west] at (axis cs:3.4,3.0){\small $\mathcal C_f$};
\end{axis}
\node[align=left,text width=3.3cm,anchor=west] at (8.6,1.6)
{\small Plus les rectangles sont fins, plus la somme de leurs aires approche $\int_a^b f$.};
\end{tikzpicture}
\fig{Figure 2 — Sommes de Riemann : l'intégrale comme limite d'aires de rectangles.}
\end{illus}

\begin{remarque}[Aire d'une fonction de signe variable]
Si $f$ \emph{change de signe} sur $[a,b]$, $\int_a^b f$ est une aire \textbf{algébrique} :
les parties où $f<0$ comptent \emph{négativement}. L'aire \emph{géométrique} (toujours positive)
du domaine entre $\mathcal C_f$ et l'axe est alors $\displaystyle\int_a^b|f(x)|\dd x$, que l'on
calcule en découpant selon le signe de $f$.
\end{remarque}

% ═══════════════════════════════════════════════════════════════
\section{Propriétés de l'intégrale}

\begin{propriete}[Premières propriétés]
Pour $f,g$ continues sur un intervalle contenant $a$ et $b$ :
\[
\int_a^a f(x)\dd x=0,\qquad \int_b^a f(x)\dd x=-\int_a^b f(x)\dd x.
\]
\end{propriete}

\begin{theoreme}[Relation de Chasles]
Pour $a,b,c$ dans $I$ (dans un ordre quelconque) :
\[
\int_a^b f(x)\dd x=\int_a^c f(x)\dd x+\int_c^b f(x)\dd x.
\]
\end{theoreme}

\begin{remarque}
Chasles sert à \textbf{découper} une intégrale là où $f$ change de forme (valeur absolue,
fonction définie par morceaux) ou de signe.
\end{remarque}

\begin{theoreme}[Linéarité]
Pour tous réels $\lambda,\mu$ :
\[
\int_a^b\big(\lambda f(x)+\mu g(x)\big)\dd x=\lambda\int_a^b f(x)\dd x+\mu\int_a^b g(x)\dd x.
\]
\end{theoreme}

\begin{exemple}
$\displaystyle\int_0^1\big(3x^2-2e^x\big)\dd x=3\int_0^1 x^2\dd x-2\int_0^1 e^x\dd x
=3\cdot\frac13-2(e-1)=1-2e+2=3-2e$.
\end{exemple}

\begin{theoreme}[Positivité et ordre]
On suppose $a\le b$.
\begin{itemize}
\item Si $f\ge0$ sur $[a,b]$, alors $\displaystyle\int_a^b f(x)\dd x\ge0$ (\textbf{positivité}).
\item Si $f\le g$ sur $[a,b]$, alors $\displaystyle\int_a^b f(x)\dd x\le\int_a^b g(x)\dd x$ (\textbf{croissance}).
\end{itemize}
\end{theoreme}

\begin{remarque}
Attention : la positivité \emph{suppose} $a\le b$. Pour $a>b$ les inégalités changent de sens.
De plus, pour $a\le b$, $\displaystyle\left|\int_a^b f\right|\le\int_a^b|f|$ (inégalité
triangulaire intégrale). Enfin, si $f$ est continue, \emph{positive} et $\int_a^b f=0$
(avec $a<b$), alors $f$ est \textbf{identiquement nulle} sur $[a,b]$.
\end{remarque}

\begin{exemple}
Pour $x\in[0,1]$, $0\le x^3\le x^2$, donc par croissance
$\displaystyle\int_0^1 x^3\dd x\le\int_0^1 x^2\dd x$, c.-à-d. $\tfrac14\le\tfrac13$. \checkmark
\end{exemple}

\begin{aretenir}
\textbf{Boîte à outils des propriétés.} Chasles (découper), linéarité (séparer/factoriser),
positivité et croissance (\emph{comparer} ou \emph{encadrer} sans calculer). Ces trois familles
suffisent à traiter la majorité des questions « sans primitive explicite ».
\end{aretenir}

% ═══════════════════════════════════════════════════════════════
\section{Valeur moyenne et inégalité de la moyenne}

\begin{definition}[Valeur moyenne]
La \textbf{valeur moyenne} de $f$ continue sur $[a,b]$ (avec $a<b$) est
\[
\mu=\frac{1}{b-a}\int_a^b f(x)\dd x.
\]
\end{definition}

\begin{remarque}
Géométriquement, $\mu$ est la hauteur du rectangle de base $[a,b]$ ayant la \emph{même aire}
que le domaine sous $\mathcal C_f$ : $\mu\,(b-a)=\int_a^b f$.
\end{remarque}

\begin{illus}
\begin{tikzpicture}[>=Stealth]
\begin{axis}[width=8.4cm,height=5cm,axis lines=middle,xlabel={\small $x$},ylabel={\small $y$},
  xmin=-0.3,xmax=4,ymin=-0.3,ymax=3.4,xtick={0.6,3.2},xticklabels={$a$,$b$},ytick={1.55},
  yticklabels={$\mu$},clip=false]
\addplot[onbsoft,domain=0.6:3.2,samples=60,fill=onbsoft,draw=none]{0.35*x^2-0.2*x+0.9}\closedcycle;
\addplot[onbo,line width=1pt,domain=0.2:3.6,samples=70]{0.35*x^2-0.2*x+0.9};
\addplot[onbblue,line width=0.9pt,domain=0.6:3.2]{1.55};
\draw[onbblue,dashed,line width=0.5pt](axis cs:0.6,0)--(axis cs:0.6,1.55);
\draw[onbblue,dashed,line width=0.5pt](axis cs:3.2,0)--(axis cs:3.2,1.55);
\node[onbo,anchor=west] at (axis cs:3.0,3.0){\small $\mathcal C_f$};
\node[onbblue,anchor=south] at (axis cs:1.9,1.55){\small $y=\mu$};
\end{axis}
\node[align=left,text width=3.2cm,anchor=west] at (8.6,1.6)
{\small Le rectangle de hauteur $\mu$ a la même aire que le domaine sous $\mathcal C_f$.};
\end{tikzpicture}
\fig{Figure 3 — Valeur moyenne : hauteur du rectangle d'aire équivalente.}
\end{illus}

\begin{theoreme}[Inégalité de la moyenne]
Si $m\le f(x)\le M$ pour tout $x\in[a,b]$ (avec $a\le b$), alors
\[
m(b-a)\le\int_a^b f(x)\dd x\le M(b-a),\qquad\text{donc}\qquad m\le\mu\le M.
\]
\end{theoreme}

\begin{theoreme}[Théorème de la moyenne]
Si $f$ est continue sur $[a,b]$ (avec $a<b$), il existe $c\in[a,b]$ tel que
$\displaystyle\int_a^b f(x)\dd x=f(c)\,(b-a)$ : la valeur moyenne est \emph{atteinte}.
\end{theoreme}

\begin{aretenir}
L'inégalité de la moyenne est l'outil-clé pour \textbf{encadrer une intégrale} qu'on ne sait
pas calculer : on majore et minore $f$ sur $[a,b]$, puis on intègre.
\end{aretenir}

% ═══════════════════════════════════════════════════════════════
\section{Intégration par parties}

\begin{theoreme}[Intégration par parties (IPP)]
Si $u$ et $v$ sont de classe $\mathcal C^1$ sur $[a,b]$ (dérivées continues), alors
\[
\int_a^b u(x)\,v'(x)\dd x=\big[u(x)\,v(x)\big]_a^b-\int_a^b u'(x)\,v(x)\dd x.
\]
\end{theoreme}

\begin{remarque}
La formule vient de $(uv)'=u'v+uv'$ : en intégrant de $a$ à $b$,
$[uv]_a^b=\int_a^b u'v+\int_a^b uv'$, d'où le résultat. On choisit $u$ que l'on \emph{dérive}
(vers quelque chose de plus simple) et $v'$ que l'on \emph{primitive}. Règle pratique de choix
de $u$ (par ordre de priorité) : \textbf{L}ogarithme, fonction \textbf{P}olynôme,
\textbf{E}xponentielle / trigonométrie (« \textbf{LPE} »). L'IPP est indispensable pour
$\int x e^{x}\dd x$, $\int x\ln x\dd x$, $\int x\cos x\dd x$, $\int\ln x\dd x$, etc.
\end{remarque}

\begin{exemple}
$\displaystyle\int_0^1 x\,e^{x}\dd x$ : posons $u=x$, $v'=e^{x}$, donc $u'=1$, $v=e^{x}$.
\[
\int_0^1 x e^{x}\dd x=\big[x e^{x}\big]_0^1-\int_0^1 e^{x}\dd x=e-\big[e^{x}\big]_0^1=e-(e-1)=1.
\]
\end{exemple}

\begin{remarque}[IPP répétée et « boucle »]
Pour $\int x^2 e^x\dd x$, on applique l'IPP \emph{deux fois} (le degré du polynôme baisse à chaque
étape). Pour $\int e^x\cos x\dd x$, deux IPP font \emph{réapparaître} l'intégrale de départ :
on la regroupe et on la résout comme une équation.
\end{remarque}

% ═══════════════════════════════════════════════════════════════
\section{Suites définies par une intégrale}

Les sujets de Bac C font très souvent intervenir une \textbf{suite} $(I_n)$ définie par une
intégrale dépendant d'un entier $n$. On y étudie : une relation de récurrence (souvent par IPP),
le signe, la monotonie, un encadrement, puis la limite.

\begin{aretenir}
\textbf{Plan d'étude d'une suite $(I_n)$ d'intégrales.}
\begin{enumerate}[label=\arabic*)]
\item \textbf{Signe} : positivité de l'intégrande $\Rightarrow$ $I_n\ge0$.
\item \textbf{Monotonie} : comparer les intégrandes de $I_{n+1}$ et $I_n$ (croissance de l'intégrale).
\item \textbf{Récurrence} : une IPP donne souvent $I_{n+1}$ en fonction de $I_n$.
\item \textbf{Encadrement} : majorer/minorer l'intégrande, puis intégrer.
\item \textbf{Limite} : théorème des gendarmes à partir de l'encadrement.
\end{enumerate}
\end{aretenir}

\begin{exemple}
$I_n=\displaystyle\int_0^1 x^n e^{x}\dd x$. Pour $x\in[0,1]$, $0\le x^n e^x\le e\,x^n$, donc
$0\le I_n\le e\displaystyle\int_0^1 x^n\dd x=\dfrac{e}{n+1}\to0$ : par les gendarmes, $\lim I_n=0$.
\end{exemple}

% ═══════════════════════════════════════════════════════════════
\section{Aires et volumes}

\begin{propriete}[Aire entre deux courbes]
Si $f$ et $g$ sont continues avec $f\le g$ sur $[a,b]$, l'aire (en u.a.) du domaine compris entre
les courbes $\mathcal C_f$ et $\mathcal C_g$ pour $x\in[a,b]$ est
\[
\mathcal A=\int_a^b\big(g(x)-f(x)\big)\dd x.
\]
\end{propriete}

\begin{illus}
\begin{tikzpicture}[>=Stealth]
\begin{axis}[width=8.6cm,height=5cm,axis lines=middle,xlabel={\small $x$},ylabel={\small $y$},
  xmin=-0.4,xmax=3.2,ymin=-0.4,ymax=3.6,xtick={0.5,2.5},xticklabels={$a$,$b$},ytick=\empty,clip=false]
% Remplissage du domaine entre les deux courbes sur [0.5 ; 2.5] :
% aire sous g (onbsoft) puis on « retire » l'aire sous f (blanc).
\addplot[draw=none,fill=onbsoft,domain=0.5:2.5,samples=60]{3-0.55*(x-1.5)^2}\closedcycle;
\addplot[draw=none,fill=white,domain=0.5:2.5,samples=60]{0.4*x+0.3}\closedcycle;
\addplot[onbo,line width=1pt,domain=0:3,samples=70]{3-0.55*(x-1.5)^2};
\addplot[onbblue,line width=1pt,domain=0:3,samples=70]{0.4*x+0.3};
\node[onbo,anchor=south] at (axis cs:1.5,2.55){\small $\mathcal C_g$};
\node[onbblue,anchor=north west] at (axis cs:2.55,1.3){\small $\mathcal C_f$};
\end{axis}
\node[align=left,text width=3.1cm,anchor=west] at (8.8,1.6)
{\small Aire entre les deux courbes : $\int_a^b\!\big(g-f\big)$, où $g$ est « au-dessus ».};
\end{tikzpicture}
\fig{Figure 4 — Aire du domaine compris entre deux courbes.}
\end{illus}

\begin{remarque}
Si les courbes se croisent, on \textbf{découpe} $[a,b]$ aux abscisses d'intersection et l'on intègre
sur chaque morceau la différence « courbe du haut $-$ courbe du bas ». L'aire géométrique totale est
$\displaystyle\int_a^b|g(x)-f(x)|\dd x$.
\end{remarque}

\begin{propriete}[Volume de révolution]
Soit $f$ continue et positive sur $[a,b]$. Le solide engendré par la rotation, autour de l'axe
$(Ox)$, du domaine sous $\mathcal C_f$ a pour volume (en unités de volume)
\[
V=\pi\int_a^b \big(f(x)\big)^2\dd x.
\]
\end{propriete}

\begin{illus}
\begin{tikzpicture}[>=Stealth]
\begin{axis}[width=8.6cm,height=4.8cm,axis lines=middle,xlabel={\small $x$},ylabel={\small $y$},
  xmin=-0.4,xmax=4.2,ymin=-2.1,ymax=2.1,xtick={0,3.5},xticklabels={$a$,$b$},ytick=\empty,clip=false]
\addplot[onbo,line width=1pt,domain=0:3.5,samples=60]{0.55*sqrt(x)+0.45};
\addplot[onbo!60,line width=0.7pt,domain=0:3.5,samples=60]{-(0.55*sqrt(x)+0.45)};
\addplot[onbsoft,draw=none,domain=0:3.5,samples=60,fill=onbsoft,opacity=0.6]{0.55*sqrt(x)+0.45}\closedcycle;
\draw[onbblue,line width=0.6pt] (axis cs:3.5,0) ellipse [x radius=0.16cm, y radius=0.95cm];
\draw[onbblue,line width=0.6pt] (axis cs:0,0) ellipse [x radius=0.08cm, y radius=0.45cm];
\end{axis}
\node[align=left,text width=3.1cm,anchor=west] at (8.8,1.6)
{\small Le domaine sous $\mathcal C_f$ tourne autour de $(Ox)$ : $V=\pi\!\int_a^b f^2$.};
\end{tikzpicture}
\fig{Figure 5 — Solide de révolution engendré par rotation autour de $(Ox)$.}
\end{illus}

\begin{exemple}
La rotation de la droite $y=x$ pour $x\in[0,h]$ autour de $(Ox)$ engendre un cône de rayon $h$ ;
$V=\pi\displaystyle\int_0^h x^2\dd x=\pi\dfrac{h^3}{3}=\dfrac{\pi h^2\cdot h}{3}=\dfrac13\pi R^2 H$
avec $R=H=h$. \checkmark\ On retrouve la formule du volume du cône.
\end{exemple}

% ═══════════════════════════════════════════════════════════════
\section{L'essentiel à retenir}

\begin{synthese}
\textbf{Définition.} $f$ continue, $F$ primitive de $f$ :\
$\displaystyle\int_a^b f(x)\dd x=\big[F(x)\big]_a^b=F(b)-F(a)$. Variable muette.\\[4pt]
\textbf{Théorème fondamental.} $\Phi(x)=\displaystyle\int_a^x f(t)\dd t$ est la primitive de $f$
nulle en $a$ : $\Phi'=f$, $\Phi(a)=0$.\\[4pt]
\textbf{Aire.} $f\ge0$ sur $[a,b]\Rightarrow\int_a^b f=$ aire sous $\mathcal C_f$ (u.a.). Aire
géométrique : $\int_a^b|f|$. Entre deux courbes ($f\le g$) : $\int_a^b(g-f)$.\\[4pt]
\textbf{Propriétés.} $\int_a^a f=0$ ;\ $\int_b^a f=-\int_a^b f$ ;\
Chasles $\int_a^b=\int_a^c+\int_c^b$ ;\ linéarité
$\int(\lambda f+\mu g)=\lambda\!\int f+\mu\!\int g$.\\[4pt]
\textbf{Ordre ($a\le b$).} $f\ge0\Rightarrow\int_a^b f\ge0$ ;\ $f\le g\Rightarrow\int_a^b f\le\int_a^b g$ ;\
$\big|\int_a^b f\big|\le\int_a^b|f|$.\\[4pt]
\textbf{Valeur moyenne.} $\mu=\dfrac1{b-a}\displaystyle\int_a^b f$. Inégalité de la moyenne :
$m\le f\le M\Rightarrow m(b-a)\le\int_a^b f\le M(b-a)$.\\[4pt]
\textbf{IPP.} $\displaystyle\int_a^b u\,v'=\big[uv\big]_a^b-\int_a^b u'\,v$. Choix de $u$ : \textbf{LPE}
(Log, Polynôme, Exp/trigo). Répéter ou « boucler » si besoin.\\[4pt]
\textbf{Formes utiles.} primitives de $u'u^n$, $\tfrac{u'}{u}$, $u'e^u$, $u'\cos u$, $u'\sin u$.\\[4pt]
\textbf{Volume de révolution.} $f\ge0$ : $V=\pi\displaystyle\int_a^b f(x)^2\dd x$ (autour de $Ox$).\\[4pt]
\textbf{Suite $(I_n)$.} signe $\to$ monotonie $\to$ récurrence (IPP) $\to$ encadrement $\to$ limite
(gendarmes).\\[4pt]
\textbf{Pièges.} oublier que la positivité exige $a\le b$ ; intégrer $|f|$ sans découper aux
changements de signe ; oublier le facteur $\pi$ ou le carré dans un volume ; mal choisir $u$ et $v'$.
\end{synthese}

% ═══════════════════════════════════════════════════════════════
\section{Méthodes \& savoir-faire}

\methode{Calculer une intégrale à l'aide d'une primitive}
Déterminer une primitive $F$ de $f$ (formes de référence : $u'u^n$, $u'/u$, $u'e^u$, $u'\sin u$\dots),
puis appliquer $\int_a^b f=F(b)-F(a)$.
\begin{exemple}
$\displaystyle\int_0^1\frac{2x}{x^2+1}\dd x$. Ici $\dfrac{2x}{x^2+1}=\dfrac{u'}{u}$ avec $u=x^2+1$, primitive $\ln u$ :
\[
\int_0^1\frac{2x}{x^2+1}\dd x=\big[\ln(x^2+1)\big]_0^1=\ln 2-\ln 1=\ln 2.
\]
\end{exemple}

\methode{Reconnaître une forme « avec $u$ »}
Repérer un facteur $u'$ devant une fonction de $u$ ; ajuster une constante multiplicative si besoin.
\begin{exemple}
$\displaystyle\int_0^1 x\,e^{x^2}\dd x$ : avec $u=x^2$, $u'=2x$, donc $x e^{x^2}=\tfrac12 u'e^u$ :
\[
\int_0^1 x e^{x^2}\dd x=\Big[\tfrac12 e^{x^2}\Big]_0^1=\tfrac12(e-1).
\]
\end{exemple}

\methode{Intégrer une fonction avec valeur absolue (ou par morceaux)}
Étudier le signe de l'expression sous la valeur absolue, puis découper l'intégrale par la relation de Chasles.
\begin{exemple}
$\displaystyle\int_{-1}^{2}|x|\dd x$ : $|x|=-x$ sur $[-1,0]$, $|x|=x$ sur $[0,2]$. Donc
\[
\int_{-1}^{2}|x|\dd x=\int_{-1}^0(-x)\dd x+\int_0^2 x\dd x=\tfrac12+2=\tfrac52.
\]
\end{exemple}

\methode{Intégration par parties (et IPP répétée)}
Choisir $u$ (à dériver) et $v'$ (à primitiver) selon la règle LPE, appliquer la formule ; recommencer
si l'intégrale restante n'est pas immédiate.
\begin{exemple}
$\displaystyle\int_1^e x\ln x\dd x$ : $u=\ln x$, $v'=x$, donc $u'=\dfrac1x$, $v=\dfrac{x^2}{2}$.
\[
\int_1^e x\ln x\dd x=\left[\frac{x^2}{2}\ln x\right]_1^e-\int_1^e\frac{x^2}{2}\cdot\frac1x\dd x
=\frac{e^2}{2}-\frac12\int_1^e x\dd x=\frac{e^2}{2}-\frac{e^2-1}{4}=\frac{e^2+1}{4}.
\]
\end{exemple}

\methode{Calculer une aire entre deux courbes}
Déterminer les abscisses d'intersection (en résolvant $f(x)=g(x)$), étudier le signe de $g-f$ pour
savoir quelle courbe est au-dessus, puis intégrer $g-f$ (en valeur absolue par morceaux si besoin).
\begin{exemple}
Aire entre $\mathcal C_g$ ($g(x)=2x$) et $\mathcal C_f$ ($f(x)=x^2$). Intersections : $x^2=2x\Rightarrow x=0$ ou $x=2$ ;
sur $[0,2]$, $2x\ge x^2$. Donc
\[
\mathcal A=\int_0^2(2x-x^2)\dd x=\left[x^2-\frac{x^3}{3}\right]_0^2=4-\frac{8}{3}=\frac{4}{3}\ \text{u.a.}
\]
\end{exemple}

\methode{Calculer un volume de révolution}
Vérifier que $f\ge0$, puis appliquer $V=\pi\displaystyle\int_a^b (f(x))^2\dd x$.
\begin{exemple}
Rotation de $\mathcal C_f$, $f(x)=\sqrt{x}$, pour $x\in[0,4]$ autour de $(Ox)$ :
\[
V=\pi\int_0^4(\sqrt x)^2\dd x=\pi\int_0^4 x\dd x=\pi\left[\frac{x^2}{2}\right]_0^4=8\pi\ \text{u.v.}
\]
\end{exemple}

\methode{Encadrer une intégrale}
Majorer et minorer $f$ sur $[a,b]$, puis intégrer (inégalité de la moyenne ou croissance).
\begin{exemple}
Pour $x\in[0,1]$ : $1\le 1+x^2\le 2$, donc $\dfrac12\le\dfrac{1}{1+x^2}\le1$, d'où
$\dfrac12\le\displaystyle\int_0^1\frac{\dd x}{1+x^2}\le1$. (Valeur exacte : $\arctan 1=\tfrac\pi4\approx0{,}785$.)
\end{exemple}

\methode{Calculer une valeur moyenne}
Calculer $\int_a^b f$, puis diviser par $b-a$.
\begin{exemple}
Valeur moyenne de $f(x)=\sin x$ sur $[0,\pi]$ :
$\mu=\dfrac1{\pi-0}\displaystyle\int_0^\pi\sin x\dd x=\dfrac1\pi\big[-\cos x\big]_0^\pi=\dfrac{2}{\pi}\approx0{,}637$.
\end{exemple}

\methode{Étudier une suite d'intégrales $I_n$}
Établir une relation de récurrence (souvent par IPP), prouver monotonie / encadrement, puis en déduire
la limite (théorème des gendarmes).
\begin{exemple}
$I_n=\displaystyle\int_0^1 x^n e^{x}\dd x$. Pour $x\in[0,1]$, $0\le x^n e^x\le e\,x^n$, donc
$0\le I_n\le e\displaystyle\int_0^1 x^n\dd x=\dfrac{e}{n+1}\to0$ : par les gendarmes, $\lim I_n=0$.
\end{exemple}

% ═══════════════════════════════════════════════════════════════
\section{Exercices d'application}

\rubrique{Calculs directs (primitives)}
\exo{}\diff{1} Calculer $\displaystyle\int_0^2(3x^2-2x+1)\dd x$, $\displaystyle\int_1^e\frac{1}{x}\dd x$
et $\displaystyle\int_0^{\pi/2}\cos x\dd x$.

\exo{}\diff{1} Calculer $\displaystyle\int_0^1 e^{2x}\dd x$ et $\displaystyle\int_0^1\frac{x}{x^2+1}\dd x$.

\exo{}\diff{2} Calculer $\displaystyle\int_0^1\frac{x}{\sqrt{x^2+1}}\dd x$ et
$\displaystyle\int_0^{\pi/2}\sin x\cos x\dd x$ (reconnaître une forme « avec $u$ »).

\exo{}\diff{2} À l'aide de la relation de Chasles, calculer $\displaystyle\int_{-1}^{2}|x|\dd x$.

\exo{}\diff{2} Calculer $\displaystyle\int_0^3|x-1|\dd x$ en découpant selon le signe de $x-1$.

\rubrique{Intégration par parties}
\exo{}\diff{2} Calculer $\displaystyle\int_0^{\pi} x\sin x\dd x$.

\exo{}\diff{2} Calculer $\displaystyle\int_1^e \ln x\dd x$ (poser $u=\ln x$, $v'=1$).

\exo{}\diff{3} En appliquant deux IPP successives, calculer $\displaystyle\int_0^1 x^2 e^{x}\dd x$.

\exo{}\diff{3} À l'aide de deux IPP, calculer $\displaystyle J=\int_0^{\pi} e^{x}\cos x\dd x$
(on fera « réapparaître » $J$).

\rubrique{Aires et volumes}
\exo{}\diff{2} Soit $f(x)=x^2-1$ et $g(x)=1-x^2$. Déterminer l'aire du domaine compris entre $\mathcal C_f$
et $\mathcal C_g$.

\exo{}\diff{2} Calculer l'aire du domaine délimité par $\mathcal C_f$ ($f(x)=4-x^2$) et l'axe $(Ox)$.

\exo{}\diff{2} Un artisan de Douala fabrique un bol obtenu en faisant tourner $\mathcal C_f$, $f(x)=\sqrt{x+1}$,
pour $x\in[0,3]$, autour de $(Ox)$. Calculer le volume intérieur (en u.v.).

\exo{}\diff{3} On considère la rotation autour de $(Ox)$ du domaine sous $\mathcal C_f$,
$f(x)=e^{-x}$, pour $x\in[0,1]$. Calculer le volume engendré.

\rubrique{Encadrement et valeur moyenne}
\exo{}\diff{2} Calculer la valeur moyenne de $f(x)=x^2$ sur $[0,3]$.

\exo{}\diff{2} Montrer que $\displaystyle 1\le\int_0^1\sqrt{1+x^3}\dd x\le\sqrt2$.

\exo{}\diff{3} Montrer que $\displaystyle\frac{1}{2}\le\int_0^1\frac{1}{1+x}\dd x\le 1$, puis calculer la valeur exacte de cette intégrale.

\rubrique{Suites d'intégrales — type Baccalauréat}
\exo{}\diff{3} On pose, pour tout entier $n\ge0$,
\[
I_n=\int_0^1 x^n e^{-x}\dd x.
\]
\begin{enumerate}[label=\arabic*)]
\item Calculer $I_0$.
\item Montrer, à l'aide d'une intégration par parties, que pour tout $n\ge0$ :
$I_{n+1}=(n+1)I_n-e^{-1}$.
\item En déduire la valeur exacte de $I_1$ et de $I_2$.
\item Montrer que pour tout $n$, $I_n\ge0$, puis que la suite $(I_n)$ est décroissante.
\item Établir l'encadrement $0\le I_n\le\dfrac{1}{n+1}$ et en déduire $\displaystyle\lim_{n\to+\infty}I_n$.
\item On considère le domaine $\mathcal D$ délimité par $\mathcal C_h$ ($h(x)=x e^{-x}$), l'axe $(Ox)$
et les droites $x=0$, $x=1$. Calculer l'aire de $\mathcal D$ et l'exprimer à l'aide de $I_1$.
\end{enumerate}

\exo{}\diff{3} Pour $n\ge1$, on pose $\displaystyle u_n=\int_0^1\frac{x^n}{1+x}\dd x$.
\begin{enumerate}[label=\arabic*)]
\item Montrer que $(u_n)$ est positive et décroissante.
\item Montrer que pour tout $n\ge1$, $\displaystyle u_n+u_{n+1}=\frac1{n+1}$.
\item En déduire un encadrement de $u_n$ puis $\displaystyle\lim_{n\to+\infty}u_n$.
\end{enumerate}

% ═══════════════════════════════════════════════════════════════
\section{Exercices d'entraînement}

\rubrique{Calculs directs et primitives}
\exo{}\diff{1} Calculer $\displaystyle\int_{-1}^{1}(x^3-x)\dd x$ et $\displaystyle\int_1^2\frac{1}{x^2}\dd x$.

\exo{}\diff{1} Calculer $\displaystyle\int_0^{\pi}\sin x\dd x$ et $\displaystyle\int_0^{\pi/4}\frac{\dd x}{\cos^2 x}$.

\exo{}\diff{1} Calculer $\displaystyle\int_0^1(2x+1)^4\dd x$ et $\displaystyle\int_0^2 e^{-x}\dd x$.

\exo{}\diff{2} Calculer $\displaystyle\int_1^4\frac{1}{\sqrt x}\dd x$ et $\displaystyle\int_0^1\frac{x^2}{x^3+1}\dd x$.

\exo{}\diff{2} Calculer $\displaystyle\int_0^{\pi/2}\cos x\,e^{\sin x}\dd x$.

\exo{}\diff{2} Calculer $\displaystyle\int_2^3\frac{2x-1}{x^2-x}\dd x$.

\exo{}\diff{2} Calculer $\displaystyle\int_{-2}^{2}|x^2-1|\dd x$.

\rubrique{Intégration par parties}
\exo{}\diff{2} Calculer $\displaystyle\int_0^1 x\,e^{2x}\dd x$.

\exo{}\diff{2} Calculer $\displaystyle\int_1^e x^2\ln x\dd x$.

\exo{}\diff{2} Calculer $\displaystyle\int_0^{\pi/2} x\cos x\dd x$.

\exo{}\diff{3} Calculer $\displaystyle\int_0^1 x^2 e^{-x}\dd x$ (deux IPP).

\exo{}\diff{3} Calculer $\displaystyle\int_0^{\pi} e^{x}\sin x\dd x$ (faire réapparaître l'intégrale).

\exo{}\diff{3} Calculer $\displaystyle\int_1^e (\ln x)^2\dd x$.

\rubrique{Aires et volumes}
\exo{}\diff{2} Calculer l'aire du domaine délimité par $\mathcal C_f$ ($f(x)=x^2$) et la droite $y=x$.

\exo{}\diff{2} Calculer l'aire entre $\mathcal C_f$ ($f(x)=x^3$) et $\mathcal C_g$ ($g(x)=x$) sur $[-1,1]$.

\exo{}\diff{2} Soit $f(x)=\sin x$ sur $[0,\pi]$. Calculer l'aire sous $\mathcal C_f$, puis le volume
engendré par rotation autour de $(Ox)$ (on rappelle $\sin^2 x=\frac{1-\cos 2x}{2}$).

\exo{}\diff{2} Une borne fontaine de Yaoundé a la forme du solide obtenu en faisant tourner
$\mathcal C_f$, $f(x)=\sqrt{4-x}$ pour $x\in[0,4]$, autour de $(Ox)$. Calculer son volume.

\exo{}\diff{3} Calculer l'aire du domaine compris entre $\mathcal C_f$ ($f(x)=\dfrac1x$), l'axe $(Ox)$
et les droites $x=1$, $x=3$.

\exo{}\diff{3} Le domaine sous $\mathcal C_f$ ($f(x)=x e^{-x}$) sur $[0,2]$ tourne autour de $(Ox)$.
Exprimer le volume engendré comme une intégrale, puis la calculer par IPP.

\rubrique{Valeur moyenne, encadrements, comparaisons}
\exo{}\diff{1} Calculer la valeur moyenne de $f(x)=2x+1$ sur $[0,4]$.

\exo{}\diff{2} Calculer la valeur moyenne de $f(x)=e^{x}$ sur $[0,1]$.

\exo{}\diff{2} Sans calculer l'intégrale, montrer que $\displaystyle 0\le\int_0^1\frac{x}{1+x^2}\dd x\le\frac12$.

\exo{}\diff{2} Montrer que $\displaystyle\frac{1}{e}\le\int_0^1 e^{-x^2}\dd x\le1$.

\exo{}\diff{2} Comparer, sans les calculer, $\displaystyle\int_0^1 x^2\,e^{x}\dd x$ et
$\displaystyle\int_0^1 x^3\,e^{x}\dd x$.

\exo{}\diff{3} Montrer que $\displaystyle\frac{\pi}{6}\le\int_0^{1/2}\frac{\dd x}{\sqrt{1-x^2}}$ est faux,
mais que $\displaystyle\int_0^{1/2}\frac{\dd x}{\sqrt{1-x^2}}\le\frac{1}{2}\cdot\frac{2}{\sqrt3}$.

\rubrique{Suites d'intégrales}
\exo{}\diff{2} Pour $n\ge0$, $\displaystyle I_n=\int_0^1 x^n\dd x$. Calculer $I_n$ et $\lim I_n$.

\exo{}\diff{3} Pour $n\ge0$, $\displaystyle J_n=\int_0^1 (1-x)^n e^x\dd x$. Établir une relation entre
$J_{n+1}$ et $J_n$, puis étudier $\lim J_n$.

\exo{}\diff{3} Pour $n\ge1$, $\displaystyle W_n=\int_0^{\pi/2}\sin^n x\dd x$ (intégrales de Wallis).
Montrer que $W_{n+2}=\dfrac{n+1}{n+2}W_n$ ; calculer $W_0,W_1,W_2,W_3$.

\exo{}\diff{3} Pour $n\ge0$, $\displaystyle K_n=\int_0^1 \frac{x^n}{1+x^2}\dd x$. Montrer que
$(K_n)$ décroît, que $K_{n+2}+K_n=\dfrac1{n+1}$, et déterminer $\lim K_n$.

\rubrique{Fonctions définies par une intégrale}
\exo{}\diff{2} Soit $\displaystyle F(x)=\int_1^x\frac{\dd t}{t}$ pour $x>0$. Calculer $F'(x)$ et
identifier $F$.

\exo{}\diff{3} Soit $\displaystyle G(x)=\int_0^x (t-1)e^{t}\dd t$. Calculer $G'(x)$, étudier les
variations de $G$, puis donner l'expression explicite de $G(x)$.

\exo{}\diff{3} Soit $\displaystyle H(x)=\int_x^{2x}\frac{\dd t}{t}$ pour $x>0$. Calculer $H(x)$
explicitement et en déduire $H'(x)$ ; commenter.

% ═══════════════════════════════════════════════════════════════
\section{Sujets type Baccalauréat}

\sujetbac[suite d'intégrales, encadrement et aire]
Pour tout entier $n\ge0$, on pose
\[
I_n=\int_0^1 \frac{x^n}{1+x}\dd x.
\]
\begin{enumerate}[label=\arabic*)]
\item Calculer $I_0$.
\item Montrer que pour tout $n\ge0$, $I_n+I_{n+1}=\dfrac{1}{n+1}$.
\item En déduire les valeurs exactes de $I_1$ et $I_2$.
\item Montrer que la suite $(I_n)$ est positive et décroissante.
\item Établir que pour tout $n\ge0$, $0\le I_n\le\dfrac{1}{n+1}$. En déduire $\displaystyle\lim_{n\to+\infty}I_n$.
\item À l'aide de la relation du 2), montrer que $\displaystyle\lim_{n\to+\infty}\big(I_n+I_{n+1}\big)=0$
et retrouver le résultat précédent.
\end{enumerate}

\sujetbac[étude de fonction, aire et volume]
On considère la fonction $f$ définie sur $[0,+\infty[$ par $f(x)=x\,e^{-x}$, de courbe $\mathcal C_f$.
\begin{enumerate}[label=\arabic*)]
\item Étudier les variations de $f$ sur $[0,+\infty[$ et préciser son maximum.
\item Déterminer $\displaystyle\lim_{x\to+\infty}f(x)$ ; en déduire une asymptote.
\item À l'aide d'une intégration par parties, calculer $\displaystyle A(t)=\int_0^t x\,e^{-x}\dd x$
pour $t>0$.
\item En déduire l'aire $\mathcal A$ du domaine compris entre $\mathcal C_f$, l'axe $(Ox)$ et les
droites $x=0$ et $x=1$.
\item Calculer $\displaystyle\lim_{t\to+\infty}A(t)$ et interpréter géométriquement.
\end{enumerate}

\sujetbac[logarithme, IPP et valeur moyenne]
On pose $f(x)=\ln x$ sur $]0,+\infty[$ et l'on s'intéresse au domaine $\mathcal D$ délimité par
$\mathcal C_f$, l'axe $(Ox)$ et la droite $x=e$.
\begin{enumerate}[label=\arabic*)]
\item Déterminer le signe de $f$ sur $[1,e]$, puis sur $[0,1]$.
\item À l'aide d'une intégration par parties, calculer $\displaystyle\int_1^e \ln x\dd x$.
\item Calculer l'aire de la partie de $\mathcal D$ située entre $x=1$ et $x=e$.
\item Calculer la valeur moyenne de $f$ sur $[1,e]$.
\item Calculer le volume du solide engendré par la rotation, autour de $(Ox)$, du domaine sous
$\mathcal C_g$ où $g(x)=\sqrt{\ln x}$ sur $[1,e]$.
\end{enumerate}

% ═══════════════════════════════════════════════════════════════
\section{Corrigés}

\corrige{de l'exercice 1}
$\displaystyle\int_0^2(3x^2-2x+1)\dd x=\big[x^3-x^2+x\big]_0^2=8-4+2=6$.\;
$\displaystyle\int_1^e\frac{\dd x}{x}=\big[\ln x\big]_1^e=1$.\;
$\displaystyle\int_0^{\pi/2}\cos x\dd x=\big[\sin x\big]_0^{\pi/2}=1$.

\corrige{de l'exercice 2}
$\displaystyle\int_0^1 e^{2x}\dd x=\left[\frac12 e^{2x}\right]_0^1=\frac{e^2-1}{2}$.\;
$\displaystyle\int_0^1\frac{x}{x^2+1}\dd x=\left[\frac12\ln(x^2+1)\right]_0^1=\frac{\ln 2}{2}$.

\corrige{de l'exercice 3}
Forme $u'/(2\sqrt u)$ avec $u=x^2+1$ ($u'=2x$) :
$\displaystyle\int_0^1\frac{x}{\sqrt{x^2+1}}\dd x=\big[\sqrt{x^2+1}\big]_0^1=\sqrt2-1$.\;
Avec $u=\sin x$ ($u'=\cos x$) : $\displaystyle\int_0^{\pi/2}\sin x\cos x\dd x
=\left[\frac{\sin^2 x}{2}\right]_0^{\pi/2}=\frac12$.

\corrige{de l'exercice 4}
Sur $[-1,0]$, $|x|=-x$ ; sur $[0,2]$, $|x|=x$. Par Chasles :
\[
\int_{-1}^{2}|x|\dd x=\int_{-1}^0(-x)\dd x+\int_0^2 x\dd x=\left[-\frac{x^2}{2}\right]_{-1}^0+\left[\frac{x^2}{2}\right]_0^2=\frac12+2=\frac52.
\]

\corrige{de l'exercice 5}
$x-1\le0$ sur $[0,1]$ ($|x-1|=1-x$), $x-1\ge0$ sur $[1,3]$ ($|x-1|=x-1$) :
\[
\int_0^3|x-1|\dd x=\int_0^1(1-x)\dd x+\int_1^3(x-1)\dd x
=\left[x-\frac{x^2}{2}\right]_0^1+\left[\frac{x^2}{2}-x\right]_1^3=\frac12+2=\frac52.
\]

\corrige{de l'exercice 6}
$u=x$, $v'=\sin x$, $u'=1$, $v=-\cos x$ :
\[
\int_0^{\pi}x\sin x\dd x=\big[-x\cos x\big]_0^{\pi}+\int_0^{\pi}\cos x\dd x=\pi+\big[\sin x\big]_0^{\pi}=\pi+0=\pi.
\]

\corrige{de l'exercice 7}
$u=\ln x$, $v'=1$, $u'=\dfrac1x$, $v=x$ :
\[
\int_1^e\ln x\dd x=\big[x\ln x\big]_1^e-\int_1^e x\cdot\frac1x\dd x=e-\big[x\big]_1^e=e-(e-1)=1.
\]

\corrige{de l'exercice 8}
$u=x^2$, $v'=e^x$ : $\displaystyle\int_0^1 x^2 e^x\dd x=\big[x^2 e^x\big]_0^1-2\int_0^1 x e^x\dd x=e-2\int_0^1 xe^x\dd x$.
Or (vu en cours) $\displaystyle\int_0^1 x e^x\dd x=1$. Donc $\displaystyle\int_0^1 x^2 e^x\dd x=e-2$.

\corrige{de l'exercice 9}
Soit $J=\displaystyle\int_0^{\pi} e^x\cos x\dd x$. IPP $u=\cos x$, $v'=e^x$ :
$J=\big[e^x\cos x\big]_0^{\pi}+\int_0^{\pi}e^x\sin x\dd x=(-e^\pi-1)+\int_0^\pi e^x\sin x\dd x$.
Seconde IPP sur $\int e^x\sin x$ ($u=\sin x$, $v'=e^x$) :
$\int_0^\pi e^x\sin x\dd x=\big[e^x\sin x\big]_0^\pi-\int_0^\pi e^x\cos x\dd x=0-J=-J$.
Donc $J=(-e^\pi-1)-J$, soit $2J=-e^\pi-1$ et $\boxed{J=-\dfrac{e^\pi+1}{2}}$.

\corrige{de l'exercice 10}
$g(x)-f(x)=(1-x^2)-(x^2-1)=2-2x^2=2(1-x^2)\ge0$ sur $[-1,1]$, et $\mathcal C_f,\mathcal C_g$ se coupent en $x=\pm1$.
\[
\mathcal A=\int_{-1}^1 2(1-x^2)\dd x=2\left[x-\frac{x^3}{3}\right]_{-1}^1=2\left(\frac23-\Big(-\frac23\Big)\right)=2\cdot\frac43=\frac83\ \text{u.a.}
\]

\corrige{de l'exercice 11}
$f(x)=4-x^2\ge0$ sur $[-2,2]$ (et $f=0$ en $\pm2$). L'aire est
\[
\int_{-2}^2(4-x^2)\dd x=\left[4x-\frac{x^3}{3}\right]_{-2}^2=\Big(8-\tfrac83\Big)-\Big(-8+\tfrac83\Big)=\frac{32}{3}\ \text{u.a.}
\]

\corrige{de l'exercice 12}
$f\ge0$ sur $[0,3]$ et $(f(x))^2=x+1$ :
\[
V=\pi\int_0^3(x+1)\dd x=\pi\left[\frac{x^2}{2}+x\right]_0^3=\pi\left(\frac92+3\right)=\frac{15\pi}{2}\ \text{u.v.}
\]

\corrige{de l'exercice 13}
$f(x)=e^{-x}\ge0$, $(f(x))^2=e^{-2x}$ :
\[
V=\pi\int_0^1 e^{-2x}\dd x=\pi\left[-\frac12 e^{-2x}\right]_0^1=\pi\cdot\frac{1-e^{-2}}{2}=\frac{\pi(1-e^{-2})}{2}\ \text{u.v.}
\]

\corrige{de l'exercice 14}
$\mu=\dfrac{1}{3-0}\displaystyle\int_0^3 x^2\dd x=\dfrac13\left[\dfrac{x^3}{3}\right]_0^3=\dfrac13\cdot 9=3$.

\corrige{de l'exercice 15}
Pour $x\in[0,1]$ : $0\le x^3\le1$ donc $1\le 1+x^3\le2$, d'où $1\le\sqrt{1+x^3}\le\sqrt2$. En intégrant
sur $[0,1]$ (longueur $1$) : $\displaystyle 1\le\int_0^1\sqrt{1+x^3}\dd x\le\sqrt2$.

\corrige{de l'exercice 16}
Pour $x\in[0,1]$ : $1\le 1+x\le2$, donc $\dfrac12\le\dfrac{1}{1+x}\le1$. En intégrant sur $[0,1]$ (longueur $1$) :
$\dfrac12\le\displaystyle\int_0^1\frac{\dd x}{1+x}\le1$. Valeur exacte : $\displaystyle\int_0^1\frac{\dd x}{1+x}=\big[\ln(1+x)\big]_0^1=\ln 2\approx0{,}693$, qui vérifie bien l'encadrement.

\corrige{de l'exercice 17}
\textbf{1)} $I_0=\displaystyle\int_0^1 e^{-x}\dd x=\big[-e^{-x}\big]_0^1=-e^{-1}+1=1-e^{-1}$.

\noindent\textbf{2)} IPP avec $u=x^{n+1}$, $v'=e^{-x}$, donc $u'=(n+1)x^{n}$, $v=-e^{-x}$ :
\[
I_{n+1}=\int_0^1 x^{n+1}e^{-x}\dd x=\big[-x^{n+1}e^{-x}\big]_0^1+\int_0^1(n+1)x^{n}e^{-x}\dd x=-e^{-1}+(n+1)I_n.
\]
D'où $I_{n+1}=(n+1)I_n-e^{-1}$.

\noindent\textbf{3)} $I_1=1\cdot I_0-e^{-1}=(1-e^{-1})-e^{-1}=1-2e^{-1}$.\;
$I_2=2I_1-e^{-1}=2(1-2e^{-1})-e^{-1}=2-5e^{-1}$.

\noindent\textbf{4)} Sur $[0,1]$, $x^n e^{-x}\ge0$ donc par positivité $I_n\ge0$. De plus, pour $x\in[0,1]$,
$x^{n+1}\le x^n$ (car $0\le x\le1$), donc $x^{n+1}e^{-x}\le x^n e^{-x}$, et par croissance de l'intégrale
$I_{n+1}\le I_n$ : la suite $(I_n)$ est décroissante.

\noindent\textbf{5)} Pour $x\in[0,1]$, $0\le e^{-x}\le1$ donc $0\le x^n e^{-x}\le x^n$ ; en intégrant :
$0\le I_n\le\displaystyle\int_0^1 x^n\dd x=\left[\frac{x^{n+1}}{n+1}\right]_0^1=\frac{1}{n+1}$.
Comme $\dfrac{1}{n+1}\to0$, le théorème des gendarmes donne $\displaystyle\lim_{n\to+\infty}I_n=0$.

\noindent\textbf{6)} Sur $[0,1]$, $h(x)=x e^{-x}\ge0$, donc l'aire de $\mathcal D$ est
\[
\mathcal A=\int_0^1 x e^{-x}\dd x=I_1=1-2e^{-1}\ \text{u.a.}\approx0{,}264\ \text{u.a.}
\]

\corrige{de l'exercice 18}
\textbf{1)} Sur $[0,1]$, $\dfrac{x^n}{1+x}\ge0$ donc $u_n\ge0$. Comme $0\le x\le1$, $x^{n+1}\le x^n$,
donc $\dfrac{x^{n+1}}{1+x}\le\dfrac{x^n}{1+x}$ et par croissance $u_{n+1}\le u_n$ : $(u_n)$ décroît.

\noindent\textbf{2)} $u_n+u_{n+1}=\displaystyle\int_0^1\frac{x^n+x^{n+1}}{1+x}\dd x=\int_0^1\frac{x^n(1+x)}{1+x}\dd x
=\int_0^1 x^n\dd x=\frac1{n+1}$.

\noindent\textbf{3)} $(u_n)$ décroît, donc $u_{n+1}\le u_n$ ; comme $u_n+u_{n+1}=\frac1{n+1}$,
on a $2u_{n+1}\le\frac1{n+1}\le 2u_n$, d'où $0\le u_n\le\frac1{n+1}\to0$. Par les gendarmes,
$\displaystyle\lim_{n\to+\infty}u_n=0$.

\corrige{du sujet 1}
\textbf{1)} $I_0=\displaystyle\int_0^1\frac{1}{1+x}\dd x=\big[\ln(1+x)\big]_0^1=\ln 2$.\\
\textbf{2)} $I_n+I_{n+1}=\displaystyle\int_0^1\frac{x^n+x^{n+1}}{1+x}\dd x=\int_0^1\frac{x^n(1+x)}{1+x}\dd x
=\int_0^1 x^n\dd x=\dfrac{1}{n+1}$.\\
\textbf{3)} $I_1=\dfrac{1}{1}-I_0=1-\ln 2$ (avec $n=0$ : $I_0+I_1=1$).\;
Avec $n=1$ : $I_1+I_2=\dfrac12$, donc $I_2=\dfrac12-(1-\ln2)=\ln 2-\dfrac12$.\\
\textbf{4)} Sur $[0,1]$, $\dfrac{x^n}{1+x}\ge0$ donc $I_n\ge0$. Et $x^{n+1}\le x^n$ sur $[0,1]$
$\Rightarrow I_{n+1}\le I_n$ : la suite décroît.\\
\textbf{5)} Pour $x\in[0,1]$, $1\le 1+x\le2$ donc $\dfrac{x^n}{2}\le\dfrac{x^n}{1+x}\le x^n$. En intégrant :
$\dfrac{1}{2(n+1)}\le I_n\le\dfrac1{n+1}$. Comme $\dfrac1{n+1}\to0$, par les gendarmes $\displaystyle\lim I_n=0$.\\
\textbf{6)} $I_n+I_{n+1}=\dfrac1{n+1}\to0$ ; comme $I_n\ge0$ et $I_{n+1}\ge0$, chacun est majoré par
$\dfrac1{n+1}\to0$, donc $\lim I_n=0$ (cohérent avec le 5).

\corrige{du sujet 2}
\textbf{1)} $f(x)=xe^{-x}$, $f'(x)=e^{-x}(1-x)$. $f'>0$ sur $[0,1[$, $f'<0$ sur $]1,+\infty[$ :
$f$ croît sur $[0,1]$, décroît sur $[1,+\infty[$. Maximum $f(1)=e^{-1}$.\\
\textbf{2)} $\displaystyle\lim_{x\to+\infty}xe^{-x}=0$ (croissances comparées) : l'axe $(Ox)$
($y=0$) est asymptote horizontale.\\
\textbf{3)} IPP $u=x$, $v'=e^{-x}$, $u'=1$, $v=-e^{-x}$ :
$A(t)=\big[-xe^{-x}\big]_0^t+\int_0^t e^{-x}\dd x=-te^{-t}+\big[-e^{-x}\big]_0^t=-te^{-t}-e^{-t}+1
=1-(t+1)e^{-t}$.\\
\textbf{4)} $f\ge0$ sur $[0,1]$, donc $\mathcal A=A(1)=1-2e^{-1}\approx0{,}264\ \text{u.a.}$\\
\textbf{5)} $\displaystyle\lim_{t\to+\infty}A(t)=\lim_{t\to+\infty}\big(1-(t+1)e^{-t}\big)=1$
(car $(t+1)e^{-t}\to0$) : l'aire sous $\mathcal C_f$ sur $[0,+\infty[$ vaut $1$ u.a.

\corrige{du sujet 3}
\textbf{1)} $\ln x\ge0\iff x\ge1$ : $f\ge0$ sur $[1,e]$ ; $f\le0$ sur $]0,1]$.\\
\textbf{2)} IPP $u=\ln x$, $v'=1$, $u'=\frac1x$, $v=x$ :
$\displaystyle\int_1^e\ln x\dd x=\big[x\ln x\big]_1^e-\int_1^e 1\dd x=e-(e-1)=1$.\\
\textbf{3)} Sur $[1,e]$, $f\ge0$ donc l'aire vaut $\displaystyle\int_1^e\ln x\dd x=1\ \text{u.a.}$\\
\textbf{4)} $\mu=\dfrac{1}{e-1}\displaystyle\int_1^e\ln x\dd x=\dfrac{1}{e-1}\approx0{,}582$.\\
\textbf{5)} $g(x)=\sqrt{\ln x}\ge0$ sur $[1,e]$ et $(g(x))^2=\ln x$ :
$V=\pi\displaystyle\int_1^e\ln x\dd x=\pi\cdot1=\pi\ \text{u.v.}$
